\section{Evaluation setup}\label{sec:evaluation}

We ask four questions.
\textbf{RQ1} Do the frozen rules certify more held-out instances than methods that also produce one policy for a
family? \textbf{RQ2} What makes them transfer: the prompt, or the verifier's refinement?
\textbf{RQ3} What do they cost, and how large are they? \textbf{RQ4} Do they hold on instances larger than any seen
in training?

\subsection{Case studies}
\paragraph*{Gridworld}
An agent visits three goals in order on an $n\times n$ grid with static obstacles and an obstacle that patrols a fixed
path. Moves succeed with probability 0.7 and slip to either side with 0.15 each. Requirements: each goal is reached
with probability at least 0.8; goals are visited in order (each ``earlier goals before goal $k$'' requirement and
their conjunction, at least 0.8); while heading to each goal, the agent never meets the moving obstacle (at least
0.7). Nine requirements per instance, all checked at once.
\emph{Random grids:} 20 grids of sizes $9$--$13$ (4 per size) from a generator that only guarantees solvability;
grids 0--11 ($9$--$11$) are for training, in four sets of a $9\times 9$, a $10\times 10$ and an $11\times 11$ grid;
grids 12--19 ($12$--$13$) are the validation set we iterated on. The sealed test set (never used before the final
runs) has 30 further grids of sizes $9$--$13$ and 12 of sizes 15, 17 and 20.
\emph{Families:} three parameterised layouts whose structure scales with $n$: \emph{rooms} (walls with one-cell doors,
a patrol in front of a door), \emph{crossing} (a wall across the grid with a single gap and a patrol along it) and
\emph{adjacent} (goals 1 and 3 side by side, so leaving goal 1 can slip into goal 3). Each has 4 members per size from
6 to 20 (seeds vary scattered obstacles). We train on sizes 6--8 (seed 0), validate on 10--14 and test on 15--20.
Every sampled member's thresholds are jointly achievable (PRISM multi-objective query).

\paragraph*{Pac-Man}
The Pac-Man MDP of the Quantitative Verification Benchmark Set~\cite{hartmanns2019qcomp}, unchanged except for the
horizon (number of moves): Pac-Man decides at eight crossings while two ghosts move randomly. Requirement: the
crash probability is at most 0.56 (the optimum is 0.5511 at every horizon). Instances are horizons 8--25; the state
space grows from \todo{k} states at 8 to 3.7 million at 25. 1--2--3--Go! reports its tree's value at horizon 25.

\subsection{Methods}
\begin{itemize}
  \item \textbf{Ours}: Claude Haiku~5.5 (and Sonnet~5.5 for one repeat), thinking enabled, 5 rounds, no restarts, the
        instruction that the rules will be certified on unseen and larger instances, and verifier refinement. Five
        repeats per condition.
  \item \textbf{1--2--3--Go!}~\cite{azeem2025onetwothreego}, reimplemented: optimal strategies of the training
        instances for one reachability objective (every goal in order, no early goal, no collision), a decision tree
        (CART, Gini, grown until pure) learned from them, applied unchanged to each target.
  \item \textbf{1--2--3--Go! with our features}: the same, with our features instead of raw coordinates as inputs.
  \item \textbf{Rule-list search}: beam search (width 4, up to 6 rules) over conjunctions of the same features with
        fixed action orders, scored on the training instances.
  \item \textbf{PPO}~\cite{schulman2017ppo}: a policy network trained on the training instances (rewards for goals in
        order, penalties for early goals and collisions), with our features or a MiniGrid-style local view as input.
  \item \textbf{Template}: a hand-written controller over relative features (towards the next goal, avoiding blocked
        cells and, in its cautious variant, moves that can slip onto a hazard).
\end{itemize}
Every method's frozen policy is written as rules and certified by the same PRISM pipeline, so all numbers are
exact certification verdicts, not simulations.
